On considère le circuit de la
Fig.~\ref{fig:2bac_svt_national_2024_rattrapage_phys_ex2_3} tel que~:
$E=\qty{10}{\V}$ et $C=\qty{3}{\uF}$.

La résistance du générateur et celle de la bobine sont supposées négligeables.

On place l'interrupteur $\mathrm{K}$ en position~(1). Lorsque le condensateur
est totalement chargé, on bascule $\mathrm{K}$ en position~(2) à l'instant
$t_{0}=0$.

La Fig.~\ref{fig:2bac_svt_national_2024_rattrapage_phys_ex2_4} donne les
variations de l'une des formes d'énergie~: l'énergie électrique
$\mathcal{E}_{e}$ emmagasinée dans le condensateur, l'énergie magnétique
$\mathcal{E}_{m}$ emmagasinée dans la bobine ou l'énergie totale $\mathcal{E}$
du circuit.

\begin{minipage}{0.48\linewidth}
    \begin{figure}[H]
        \centering
        \begin{circuitikz}
            \newdimen\d \d=3cm
            \draw (0,0) node[cute spdt mid,rotate=90] (Sw) {};
            \node[above] at (Sw.out 1) {(1)};
            \node[above] at (Sw.out 2) {(2)};
            \node[above=9mm] at (Sw.in) {$\mathrm{K}$};
            \draw (Sw.out 1)
                 to[short] ++(-.8\d,0)
                 to[vsource,v_<=$E$,voltage shift=0.5] ++(0,-\d)
                 to[short] ++(2\d,0) coordinate(M)
                 to[L,l_={$L$},mirror,voltage shift=3.5,name=C] ++(0,\d)
                 to[short] (Sw.out 2);
            \draw (Sw.in) to[C,l_=$C$,i>^=$i$] (Sw.in |- M);
        \end{circuitikz}
    \caption{}
    \label{fig:2bac_svt_national_2024_rattrapage_phys_ex2_3}
    \end{figure}
\end{minipage}
\hfill
\begin{minipage}{0.48\linewidth}
    \begin{figure}[H]
        \centering
        \includegraphics[width=.7\textwidth]{2bac_svt_national_2024_rattrapage_phys_ex2_4}
        \caption{}
        \label{fig:2bac_svt_national_2024_rattrapage_phys_ex2_4}
    \end{figure}
\end{minipage}

\begin{enumerate}
    \item Donner, en justifiant, la forme d'énergie correspondante à la courbe
          représentée dans la
          Fig.~\ref{fig:2bac_svt_national_2024_rattrapage_phys_ex2_4}.
    \item Expliquer de point de vue énergétique le régime d'oscillations
          électriques obtenu.
    \item Déterminer graphiquement la valeur de la période propre $T_{0}$ des
          oscillations.
    \item Déterminer la valeur de l'inductance $L$ de la bobine.
    \item Déterminer la valeur de l'énergie magnétique $\mathcal{E}_{m}$
          emmagasinée dans la bobine à l'instant $t=\frac{3 T_{0}}{4}$.
\end{enumerate}

